\begin{figure*}[t]
\centering
\figplaceholder[2.0in]{\textwidth}{Fig.~4 --- the fourteen \bench tasks}{%
\textbf{Layout:} full width, two rows of seven equal panels, one per task, rendered agent-view frames of \bench (no photographs). A thin family label bar spans each group of panels. Row 1: \emph{Grasp} (harmonic, conveyor, translating, rotating, occluded) and \emph{Place} (moving receptacle, rotating receptacle). Row 2: \emph{Pass} (wire passage, full wire passage), \emph{Capture} (ball, ball on table, slider, can), and \emph{Impact} (post-impact grasp). Task name under each panel.
\textbf{Each panel:} cropped to the region in which the target moves, at the same scale in all panels; 3--4 semi-transparent past poses of the target and an arrow for its motion, so that the motion pattern is visible (oscillation, straight line, rotation, disappearance behind the occluder, rolling, deflection after the impact); obstacles outlined and drivers marked in the role colors of Fig.~\ref{fig:principles}c.
\textbf{Rules:} no timeline (Figs.~\ref{fig:architecture} and~\ref{fig:principles} already show one); same camera in all panels.}
\caption{\textbf{\bench.} Fourteen tasks in five families, each shown with the recent
motion of its target. Observation, scene motion, and execution share one physical
clock, so the scene keeps moving while a policy computes.}
\label{fig:bench_overview}
\end{figure*}